\chapter{Toolkit \& Cheat Sheets: Peralatan Praktis}

\begin{insightbox}
\textit{``Bab ini adalah \highlight{arsenal praktis} untuk hubunganmu. Print, fotokopi, tempel di kulkas --- gunakan ini sebagai \highlight{reference guide} saat kamu butuh bantuan cepat di tengah konflik atau ingin memperkuat hubungan.''}
\end{insightbox}

\section{Quick Reference Cards}

\subsection{Kartu Identifikasi Cepat: Island, Wave, Anchor}

\begin{center}
\renewcommand{\arraystretch}{1.5}
\begin{tabular}{|>{\raggedright}p{2.5cm}|>{\raggedright}p{4.2cm}|>{\raggedright}p{4.2cm}|>{\raggedright\arraybackslash}p{4.2cm}|}
\hline
\rowcolor{primary!20}
\textbf{Aspek} & \textbf{Anchor (Jangkar)} & \textbf{Island (Pulau)} & \textbf{Wave (Gelombang)} \\
\hline
\textbf{Motto} & Dua lebih baik dari satu & Aku bisa sendiri & Aku nggak bisa tanpamu \\
\hline
\textbf{Saat stres} & Kolaboratif, cari solusi & Menjauh, butuh ruang & Mengejar, emosional \\
\hline
\textbf{Kebutuhan utama} & Keseimbangan dekat \& jauh & Banyak ruang sendiri & Banyak kedekatan, reassurance \\
\hline
\textbf{Ketakutan terbesar} & Kehilangan koneksi yang sudah dibangun & Kehilangan otonomi, terjebak & Ditinggal, tidak dicintai \\
\hline
\textbf{Respons konflik} & Diskusi, cari win-win & Menghindar atau menyerang balik & Emosional, dramatis, testing \\
\hline
\textbf{Cara support} & Jadi partner yang stabil & Beri ruang, jangan kejar & Kejar dengan lembut, validasi \\
\hline
\end{tabular}
\end{center}

\begin{tipbox}[Checklist Identifikasi Cepat]
\textbf{$\square$ Island jika dia:}
\begin{itemize}[itemsep=0.1cm]
    \item ``Lupa'' memberi tahu kemana perginya
    \item Saat ada masalah, menghilang atau merenung sendiri
    \item Sulit bertransisi dari aktivitas solo ke interaksi
    \item Saat dipeluk terlalu lama, mulai gelisah
    \item Jawaban standar: ``Tidak apa-apa'' / ``Aku fine''
\end{itemize}

\textbf{$\square$ Wave jika dia:}
\begin{itemize}[itemsep=0.1cm]
    \item Cemas jika tidak membalas pesan dengan cepat
    \item Butuh validasi berulang-ulang
    \item Kadang mendekat berlebihan, tiba-tiba mendorong
    \item Sering test: ``Kalau aku pergi, apakah dia mengejar?''
    \item Menggunakan kata hiperbolis: ``selalu'', ``tidak pernah''
\end{itemize}

\textbf{$\square$ Anchor jika dia:}
\begin{itemize}[itemsep=0.1cm]
    \item Mudah meminta bantuan dan menerima bantuan
    \item Tidak merasa terancam jika pasangan butuh ruang
    \item Bisa sendiri atau bersama dengan sama nyamannya
    \item Mencari solusi win-win saat konflik
    \item Bisa bilang ``aku memilihmu'' tanpa ragu
\end{itemize}
\end{tipbox}

\subsection{Four Horsemen: Warning Signs \& Antidotes}

\begin{warningbox}[The Four Horsemen --- Tanda Bahaya!]
\textbf{Identifikasi dan Antidotes}

\begin{center}
\renewcommand{\arraystretch}{1.6}
\begin{tabular}{|>{\raggedright}p{3cm}|>{\raggedright}p{5cm}|>{\raggedright\arraybackslash}p{6cm}|}
\hline
\rowcolor{accent!20}
\textbf{Horseman} & \textbf{Warning Signs} & \textbf{Antidote} \\
\hline
\textbf{Criticism} & Menyerang karakter pasangan; ``Kamu selalu...'' / ``Kamu tidak pernah...'' & \textbf{Soft Start-up}: Keluhkan perilaku spesifik tanpa menyalahkan; ``Aku merasa... ketika...'' \\
\hline
\textbf{Contempt} & Merendahkan, sarkasme, eye-rolling, sinisme, ejekan & \textbf{Build Culture of Appreciation}: Fokus pada apa yang kamu syukuri, ungkapkan admirasi \\
\hline
\textbf{Defensiveness} & Membalas serangan, membuat alasan, memutarbalikkan & \textbf{Take Responsibility}: Akui bagianmu; ``Ada benarnya juga...'' \\
\hline
\textbf{Stonewalling} & Menarik diri, shutdown, tidak merespons, flooding & \textbf{Self-Soothing}: Stop, beri tahu butuh break 20 menit, kembali saat sudah tenang \\
\hline
\end{tabular}
\end{center}
\end{warningbox}

\subsection{Bank Kalimat: Accepting Influence}

\begin{praktikbox}[Phrase Bank --- Menerima Pengaruh]
Gunakan kalimat-kalimat ini untuk menunjukkan bahwa kamu \highlight{menerima pengaruh} pasangan:

\textbf{Level 1 --- Validasi Dasar:}
\begin{itemize}[itemsep=0.15cm]
    \item[$\square$] ``Itu pendapat yang bagus.''
    \item[$\square$] ``Aku belum memikirkan dari sudut pandang itu.''
    \item[$\square$] ``Kamu punya point yang valid.''
    \item[$\square$] ``Aku mengerti maksudmu.''
\end{itemize}

\textbf{Level 2 --- Menunjukkan Pemahaman:}
\begin{itemize}[itemsep=0.15cm]
    \item[$\square$] ``Aku bisa lihat kenapa kamu merasa begitu.''
    \item[$\square$] ``Jika aku jadi kamu, aku mungkin juga akan...''
    \item[$\square$] ``Masuk akal kalau kamu berpikir seperti itu.''
    \item[$\square$] ``Aku paham apa yang kamu rasakan.''
\end{itemize}

\textbf{Level 3 --- Komitmen Bersama:}
\begin{itemize}[itemsep=0.15cm]
    \item[$\square$] ``Bagaimana kalau kita coba cara kamu dulu?''
    \item[$\square$] ``Mari kita cari solusi yang cocok untuk kita berdua.''
    \item[$\square$] ``Aku akan melakukan [action] untuk menunjukkan komitmenku.''
    \item[$\square$] ``Kita adalah tim --- apa yang terbaik untuk kita?''
\end{itemize}
\end{praktikbox}

\subsection{Bank Kalimat: Repair Attempts Berdasarkan Tingkat Severity}

\begin{praktikbox}[Repair Attempt Phrase Bank]
\textbf{Level 1 --- Ringan (Konflik kecil, salah paham):}
\begin{itemize}[itemsep=0.1cm]
    \item[$\square$] ``Maaf, aku salah.''
    \item[$\square$] ``Bolehkah aku ulang?''
    \item[$\square$] ``Aku tidak maksud begitu.''
    \item[$\square$] ``Ini bukan tentang kamu, ini tentang [situasi].''
    \item[$\square$] ``Bisa kita mulai ulang?''
\end{itemize}

\textbf{Level 2 --- Sedang (Perasaan terluka, frustrasi):}
\begin{itemize}[itemsep=0.1cm]
    \item[$\square$] ``Aku tahu aku salah. Bagaimana aku bisa memperbaikinya?''
    \item[$\square$] ``Maafkan kata-kataku tadi, itu tidak pantas.''
    \item[$\square$] ``Aku menghargai kesabaranmu.''
    \item[$\square$] ``Ini penting untukku, tapi kamu juga penting.''
    \item[$\square$] ``Aku ingin kita bisa menyelesaikan ini bersama.''
\end{itemize}

\textbf{Level 3 --- Berat (Four Horsemen muncul, flooding):}
\begin{itemize}[itemsep=0.1cm]
    \item[$\square$] ``Aku butuh istirahat sebentar, tapi aku akan kembali. Ini janji.''
    \item[$\square$] ``Maafkan aku telah [specific behavior]. Itu tidak benar.''
    \item[$\square$] ``Aku memilihmu daripada masalah ini.''
    \item[$\square$] ``Aku tidak ingin bertengkar seperti ini denganmu.''
    \item[$\square$] ``Kamu lebih penting daripada memenangkan argumen ini.''
\end{itemize}
\end{praktikbox}

\section{Templates Praktis}

\subsection{Template: State of the Union Meeting}

\begin{praktikbox}[State of the Union Meeting Template]
\textbf{Jadwal:} Setiap \rule{3cm}{0.4pt} (minggu/bulan) \\
\textbf{Durasi:} 30-45 menit \\
\textbf{Tempat:} Tempat nyaman tanpa gangguan

\vspace{0.3cm}
\textbf{STRUKTUR MEETING:}

\textit{Fase 1: Apresiasi (5 menit)}
\begin{itemize}[itemsep=0.1cm]
    \item[$\square$] Apa yang aku syukuri dari pasanganku minggu ini:
    \item[$\square$] Satu hal spesifik yang dia lakukan yang membuatku merasa dicintai:
\end{itemize}

\textit{Fase 2: Update Love Maps (10 menit)}
\begin{itemize}[itemsep=0.1cm]
    \item[$\square$] Hal penting yang terjadi di hidupku minggu ini:
    \item[$\square$] Apa yang sedang aku khawatirkan/stres:
    \item[$\square$] Apa yang sedang aku nantikan:
    \item[$\square$] Pertanyaan untuk pasangan: \rule{6cm}{0.4pt}
\end{itemize}

\textit{Fase 3: Review Konflik (15-20 menit)}
\begin{itemize}[itemsep=0.1cm]
    \item[$\square$] Apakah ada konflik yang belum terselesaikan? $\square$ Ya $\square$ Tidak
    \item[$\square$] Jika ya, apakah ini solvable atau perpetual? $\square$ Solvable $\square$ Perpetual
    \item[$\square$] Jika solvable: Diskusikan solusi konkret
    \item[$\square$] Jika perpetual: Diskusikan bagaimana dialog gridlock
\end{itemize}

\textit{Fase 4: Plan Ahead (5-10 menit)}
\begin{itemize}[itemsep=0.1cm]
    \item[$\square$] Komitmen minggu depan:
    \item[$\square$] Ritual koneksi yang akan kita lakukan:
    \item[$\square$] Apa yang perlu diperbaiki:
\end{itemize}
\end{praktikbox}

\subsection{Template: Conflict Resolution Flowchart}

\begin{center}
\begin{tikzpicture}[
    node distance=1.2cm,
    decision/.style={draw=primary, diamond, fill=lightbg, 
                     text width=2.8cm, align=center, aspect=2, font=\small},
    process/.style={draw=secondary, rectangle, fill=secondary!10,
                    text width=3.5cm, align=center, rounded corners, font=\small},
    startstop/.style={draw=accent, rectangle, fill=accent!10, rounded corners=10pt,
                      text width=3cm, align=center, font=\small\bfseries},
    arrow/.style={->, thick, primary}
]
    % Nodes
    \node[startstop] (start) {KONFLIK MULAI};
    \node[decision, below=of start] (flooded) {Apakah flooded?};
    \node[process, right=1.5cm of flooded] (break) {Gunakan Safe Word\\Istirahat 20-30 menit};
    \node[process, below=of flooded] (soft) {Gunakan Soft Start-up\\``Aku marah karena...''};
    \node[decision, below=of soft] (solvable) {Solvable atau\\Perpetual?};
    \node[process, left=1.5cm of solvable] (solve) {Problem-solving\\Cari solusi konkret};
    \node[process, below=of solvable] (manage) {Dialog Gridlock\\Pahami dream masing-masing};
    \node[process, below=of manage] (repair) {Repair Attempt\\Maaf, humor, kasih sayang};
    \node[startstop, below=of repair] (end) {KONEKSI PULIH};
    
    % Arrows
    \draw[arrow] (start) -- (flooded);
    \draw[arrow] (flooded) -- node[above, font=\scriptsize] {Ya} (break);
    \draw[arrow] (flooded) -- node[left, font=\scriptsize] {Tidak} (soft);
    \draw[arrow] (break) |- (soft);
    \draw[arrow] (soft) -- (solvable);
    \draw[arrow] (solvable) -- node[above, font=\scriptsize] {Solvable} (solve);
    \draw[arrow] (solvable) -- node[right, font=\scriptsize] {Perpetual} (manage);
    \draw[arrow] (solve) |- (repair);
    \draw[arrow] (manage) -- (repair);
    \draw[arrow] (repair) -- (end);
\end{tikzpicture}
\end{center}

\begin{tipbox}[Panduan Keputusan Cepat]
\textbf{$\square$ SAFE WORD Protocol:}
\begin{enumerate}[itemsep=0.1cm]
    \item Satu pihak menggunakan safe word
    \item Konflik \textbf{langsung stop} --- tidak ada ``tambahan kata terakhir''
    \item Pergi ke ruang terpisah
    \item Self-soothe 20-30 menit (jantung kembali normal)
    \item Kembali dan coba lagi dengan soft start-up
\end{enumerate}

\textbf{$\square$ Solvable vs Perpetual:}
\begin{itemize}[itemsep=0.1cm]
    \item \textbf{Solvable:} Masalah spesifik dengan situasi spesifik $\rightarrow$ Cari solusi konkret
    \item \textbf{Perpetual:} Perbedaan karakter/latar belakang $\rightarrow$ Dialog dengan empati
\end{itemize}
\end{tipbox}

\subsection{Template: Couple Agreement}

\begin{exercisebox}[Couple Agreement --- Kontrak Hubungan Kami]
\textbf{Kami, [Nama 1] dan [Nama 2], berkomitmen untuk:}

\vspace{0.3cm}
\textbf{KOMUNIKASI}
\begin{itemize}[itemsep=0.15cm]
    \item[$\square$] Melakukan State of the Union meeting setiap \rule{3cm}{0.4pt}
    \item[$\square$] Melakukan daily check-in minimal \rule{1cm}{0.4pt} menit per hari
    \item[$\square$] Safe word kami adalah: \rule{5cm}{0.4pt}
    \item[$\square$] Tidak menggunakan Four Horsemen dalam konflik
    \item[$\square$] Selalu melakukan repair attempt setelah konflik
\end{itemize}

\textbf{KONEKSI EMOSIONAL}
\begin{itemize}[itemsep=0.15cm]
    \item[$\square$] Update Love Maps secara rutin
    \item[$\square$] Quality time minimal \rule{1cm}{0.4pt} jam per minggu
    \item[$\square$] Date night setiap \rule{3cm}{0.4pt}
    \item[$\square$] Ritual koneksi harian: \rule{8cm}{0.4pt}
\end{itemize}

\textbf{KONFLIK}
\begin{itemize}[itemsep=0.15cm]
    \item[$\square$] Jika flooded: gunakan safe word, istirahat 20-30 menit
    \item[$\square$] Selalu gunakan soft start-up
    \item[$\square$] Terima pengaruh satu sama lain
    \item[$\square$] Jika perpetual problem: lakukan dialog gridlock
\end{itemize}

\textbf{KOMITMEN KHUSUS [Nama 1]:}
\begin{itemize}[itemsep=0.15cm]
    \item[$\square$] Aku akan \rule{10cm}{0.4pt}
    \item[$\square$] Saat aku butuh ruang, aku akan \rule{8cm}{0.4pt}
\end{itemize}

\textbf{KOMITMEN KHUSUS [Nama 2]:}
\begin{itemize}[itemsep=0.15cm]
    \item[$\square$] Aku akan \rule{10cm}{0.4pt}
    \item[$\square$] Saat aku cemas, aku akan \rule{8cm}{0.4pt}
\end{itemize}

\vspace{0.5cm}
\rule{6cm}{0.4pt} \hspace{1cm} \rule{6cm}{0.4pt}\\
[Nama 1] \hspace{4cm} [Nama 2]

\vspace{0.3cm}
Tanggal: \rule{3cm}{0.4pt}
\end{exercisebox}

\subsection{Template: Annual Relationship Review}

\begin{praktikbox}[Annual Relationship Review Template]
\textbf{Dilakukan setiap tahun (misal: anniversary)}

\vspace{0.3cm}
\textbf{BAGIAN 1: REFLEKSI TAHUN INI}

\textit{Hal-hal yang berhasil dengan baik:}
\begin{itemize}[itemsep=0.1cm]
    \item[$\square$] \\
    \item[$\square$] \\
    \item[$\square$] \\
\end{itemize}

\textit{Hal-hal yang perlu diperbaiki:}
\begin{itemize}[itemsep=0.1cm]
    \item[$\square$] \\
    \item[$\square$] \\
    \item[$\square$] \\
\end{itemize}

\textit{Momen terbaik tahun ini:} \hrulefill

\textit{Tantangan terbesar yang kita atasi:} \hrulefill

\textbf{BAGIAN 2: LOVE MAPS UPDATE}
\begin{itemize}[itemsep=0.1cm]
    \item[$\square$] Tujuan/aspirasi pasanganku saat ini:
    \item[$\square$] Beban terbesar yang dia pikul:
    \item[$\square$] Hal yang sedang dia khawatirkan:
    \item[$\square$] Apa yang membuatnya merasa paling hidup:
\end{itemize}

\textbf{BAGIAN 3: PERPETUAL PROBLEMS REVIEW}
\begin{itemize}[itemsep=0.1cm]
    \item[$\square$] Daftar perpetual problems kami:
    \begin{enumerate}
        \item 
        \item 
        \item 
    \end{enumerate}
    \item[$\square$] Bagaimana progress dialog gridlock untuk masing-masing?
    \item[$\square$] Apakah ada yang perlu approach baru?
\end{itemize}

\textbf{BAGIAN 4: GOALS TAHUN DEPAN}
\begin{itemize}[itemsep=0.1cm]
    \item[$\square$] Satu hal yang ingin kami capai bersama:
    \item[$\square$] Satu ritual baru yang ingin kami coba:
    \item[$\square$] Satu area hubungan yang ingin kami perkuat:
    \item[$\square$] Satu pengalaman baru yang ingin kami coba:
\end{itemize}
\end{praktikbox}

\section{Worksheets (Lembar Kerja)}

\subsection{Worksheet: Love Maps Questionnaire (20 Pertanyaan)}

\begin{exercisebox}[Love Maps Questionnaire]
\textbf{Petunjuk:} Jawablah pertanyaan-pertanyaan ini tentang pasanganmu. Tujuannya bukan untuk mendapat skor sempurna, tapi untuk melihat seberapa baik kamu mengenalnya.

\vspace{0.3cm}
\textbf{BAGIAN A: TENTANG KEHIDUPANNYA}

\begin{enumerate}[itemsep=0.2cm]
    \item Siapa sahabat terbaik pasanganmu saat ini? \hrulefill
    \item Siapa musuh bebuyutannya saat ini? \hrulefill
    \item Apa beban terbesar yang sedang dipikulnya? \hrulefill
    \item Apa yang paling membuatnya stres minggu ini? \hrulefill
    \item Apa yang paling membuatnya bahagia akhir-akhir ini? \hrulefill
\end{enumerate}

\textbf{BAGIAN B: TENTANG IMPIAN \& TAKUTNYA}

\begin{enumerate}[itemsep=0.2cm]
    \setcounter{enumi}{5}
    \item Apa impian terbesarnya yang belum pernah dia ceritakan? \hrulefill
    \item Apa ketakutan terbesarnya tentang masa depan? \hrulefill
    \item Apa yang ingin dia capai dalam 5 tahun ke depan? \hrulefill
    \item Apa pengalaman masa kecil yang paling berkesan baginya? \hrulefill
    \item Apa yang membuatnya merasa paling insecure? \hrulefill
\end{enumerate}

\textbf{BAGIAN C: TENTANG PREFERENSINYA}

\begin{enumerate}[itemsep=0.2cm]
    \setcounter{enumi}{10}
    \item Apa makanan favoritnya? \hrulefill
    \item Apa film/serial favoritnya saat ini? \hrulefill
    \item Apa musik yang paling dia suka? \hrulefill
    \item Apa aktivitas yang paling dia nikmati untuk melepas stres? \hrulefill
    \item Apa hadiah yang paling berarti baginya? \hrulefill
\end{enumerate}

\textbf{BAGIAN D: TENTANG HUBUNGAN}

\begin{enumerate}[itemsep=0.2cm]
    \setcounter{enumi}{15}
    \item Apa yang paling dia hargai dari hubungan kita? \hrulefill
    \item Apa yang membuatnya merasa paling dicintai? \hrulefill
    \item Apa yang paling dia butuhkan dariku saat stres? \hrulefill
    \item Apa yang dia harap aku lakukan lebih sering? \hrulefill
    \item Apa makna cinta baginya? \hrulefill
\end{enumerate}

\vspace{0.3cm}
\textbf{SKORING:}
\begin{itemize}[itemsep=0.1cm]
    \item 16-20 benar: \textbf{Excellent} --- Love Maps-mu sangat detail!
    \item 11-15 benar: \textbf{Good} --- Kamu mengenalnya dengan baik, tapi masih ada ruang untuk eksplorasi.
    \item 6-10 benar: \textbf{Fair} --- Perlu update Love Maps secara rutin.
    \item 0-5 benar: \textbf{Needs Work} --- Luangkan waktu untuk mengenalnya lebih dalam.
\end{itemize}
\end{exercisebox}

\subsection{Worksheet: Emotional Bank Account Tracker}

\begin{exercisebox}[Emotional Bank Account Tracker]
\textbf{Petunjuk:} Setiap interaksi positif adalah \textbf{deposit}, setiap interaksi negatif adalah \textbf{withdrawal}. Catat selama seminggu.

\vspace{0.3cm}
\textbf{MINGGU:} \rule{3cm}{0.4pt} \hspace{1cm} \textbf{BULAN:} \rule{3cm}{0.4pt}

\begin{center}
\renewcommand{\arraystretch}{1.4}
\begin{tabular}{|p{2cm}|p{4.5cm}|p{2cm}|p{4.5cm}|p{2cm}|}
\hline
\rowcolor{sagegreen!20}
\textbf{Hari} & \textbf{Deposit (+)} & \textbf{+Poin} & \textbf{Withdrawal (-)} & \textbf{-Poin} \\
\hline
Senin & & $\square$ 1 $\square$ 2 $\square$ 3 & & $\square$ 1 $\square$ 2 $\square$ 3 \\
\hline
Selasa & & $\square$ 1 $\square$ 2 $\square$ 3 & & $\square$ 1 $\square$ 2 $\square$ 3 \\
\hline
Rabu & & $\square$ 1 $\square$ 2 $\square$ 3 & & $\square$ 1 $\square$ 2 $\square$ 3 \\
\hline
Kamis & & $\square$ 1 $\square$ 2 $\square$ 3 & & $\square$ 1 $\square$ 2 $\square$ 3 \\
\hline
Jumat & & $\square$ 1 $\square$ 2 $\square$ 3 & & $\square$ 1 $\square$ 2 $\square$ 3 \\
\hline
Sabtu & & $\square$ 1 $\square$ 2 $\square$ 3 & & $\square$ 1 $\square$ 2 $\square$ 3 \\
\hline
Minggu & & $\square$ 1 $\square$ 2 $\square$ 3 & & $\square$ 1 $\square$ 2 $\square$ 3 \\
\hline
\end{tabular}
\end{center}

\vspace{0.3cm}
\textbf{RUMUS PERHITUNGAN:}
\begin{itemize}[itemsep=0.1cm]
    \item Deposit: $\square$ 1 (kecil), $\square$ 2 (sedang), $\square$ 3 (besar)
    \item Withdrawal: $\square$ 1 (kecil), $\square$ 2 (sedang), $\square$ 3 (besar)
\end{itemize}

\textbf{RASIO MINGGU INI:}
\begin{itemize}[itemsep=0.1cm]
    \item Total Deposit: \rule{2cm}{0.4pt}
    \item Total Withdrawal: \rule{2cm}{0.4pt}
    \item Rasio (Deposit:Withdrawal): \rule{3cm}{0.4pt}
\end{itemize}

\begin{tipbox}[Target Rasio]
\textbf{Target: 5:1 atau lebih tinggi} (5 deposit untuk setiap 1 withdrawal)

Jika rasio di bawah 3:1, hubungan berada dalam zona berbahaya.
\end{tipbox}
\end{exercisebox}

\subsection{Worksheet: Shared Meaning Assessment}

\begin{exercisebox}[Shared Meaning Assessment]
\textbf{Petunjuk:} Diskusikan bersama pasangan. Tidak ada jawaban benar atau salah --- yang penting adalah keselarasan.

\vspace{0.3cm}
\textbf{BAGIAN 1: PERAN DALAM HUBUNGAN}

\textit{Apa peran yang kita inginkan dalam hubungan?}
\begin{center}
\renewcommand{\arraystretch}{1.3}
\begin{tabular}{|p{4cm}|p{4cm}|p{5cm}|}
\hline
\rowcolor{primary!15}
\textbf{Area} & \textbf{[Nama 1]} & \textbf{[Nama 2]} \\
\hline
Pengambilan keputusan besar & & \\
\hline
Pengaturan keuangan & & \\
\hline
Tanggung jawab rumah tangga & & \\
\hline
Pengasuhan anak (jika ada) & & \\
\hline
Sosial/keluarga & & \\
\hline
\end{tabular}
\end{center}

\textbf{BAGIAN 2: NILAI \& VISI}

\textit{Seberapa selaras visi kita? (1-5, 5 = sangat selaras)}

\begin{itemize}[itemsep=0.2cm]
    \item[$\square$] Keuangan: [1] [2] [3] [4] [5]
    \item[$\square$] Karir dan ambisi: [1] [2] [3] [4] [5]
    \item[$\square$] Pola asuh: [1] [2] [3] [4] [5]
    \item[$\square$] Kehidupan sosial: [1] [2] [3] [4] [5]
    \item[$\square$] Spiritualitas/keyakinan: [1] [2] [3] [4] [5]
    \item[$\square$] Gaya hidup (kesehatan, olahraga): [1] [2] [3] [4] [5]
\end{itemize}

\textbf{BAGIAN 3: RITUAL KONEKSI}

\textit{Ritual apa yang sudah kita miliki?}
\begin{itemize}[itemsep=0.15cm]
    \item[$\square$] Morning ritual: \hrulefill
    \item[$\square$] Bedtime ritual: \hrulefill
    \item[$\square$] Greeting ritual (saat bertemu): \hrulefill
    \item[$\square$] Weekend ritual: \hrulefill
    \item[$\square$] Anniversary ritual: \hrulefill
\end{itemize}

\textit{Ritual baru yang ingin kita ciptakan:} \hrulefill
\end{exercisebox}

\subsection{Worksheet: Perpetual Problem Identification}

\begin{exercisebox}[Identifikasi Perpetual Problem]
\textbf{Petunjuk:} Identifikasi masalah yang muncul berulang kali dan selalu sulit diselesaikan.

\vspace{0.3cm}
\textbf{MASLAH 1}
\begin{itemize}[itemsep=0.15cm]
    \item Deskripsi masalah: \hrulefill
    \item Topik permukaan (apa yang kita pertengkarkan): \hrulefill
    \item Dream dalam konflik ini (di balik topik permukaan): \hrulefill
    \item Bagian [Nama 1]: \hrulefill
    \item Bagian [Nama 2]: \hrulefill
\end{itemize}

\textbf{MASLAH 2}
\begin{itemize}[itemsep=0.15cm]
    \item Deskripsi masalah: \hrulefill
    \item Topik permukaan (apa yang kita pertengkarkan): \hrulefill
    \item Dream dalam konflik ini (di balik topik permukaan): \hrulefill
    \item Bagian [Nama 1]: \hrulefill
    \item Bagian [Nama 2]: \hrulefill
\end{itemize}

\textbf{MASLAH 3}
\begin{itemize}[itemsep=0.15cm]
    \item Deskripsi masalah: \hrulefill
    \item Topik permukaan (apa yang kita pertengkarkan): \hrulefill
    \item Dream dalam konflik ini (di balik topik permukaan): \hrulefill
    \item Bagian [Nama 1]: \hrulefill
    \item Bagian [Nama 2]: \hrulefill
\end{itemize}

\begin{tipbox}[Panduan Identifikasi]
Perpetual problem biasanya:
\begin{itemize}[itemsep=0.1cm]
    \item Muncul berulang kali dengan pola yang sama
    \item Tidak pernah benar-benar terselesaikan
    \item Berkaitan dengan perbedaan karakter, latar belakang, atau nilai inti
    \item Saat dialog dengan baik, bukan untuk menyelesaikan, tapi untuk memahami
\end{itemize}
\end{tipbox}
\end{exercisebox}

\section{Emergency Protocols (Protokol Darurat)}

\subsection{Protokol: Ketika Flooding Terjadi}

\begin{warningbox}[EMERGENCY PROTOCOL --- FLOODING]
\textbf{STAGE 1: DETEKSI (Cek gejala flooding)}
\begin{itemize}[itemsep=0.1cm]
    \item[$\square$] Jantung berdebar lebih dari 100 bpm
    \item[$\square$] Napas pendek/cepat
    \item[$\square$] Otot tegang (rahang, bahu, tangan)
    \item[$\square$] Mulut kering
    \item[$\square$] Pikiran berkecamuk atau blank
    \item[$\square$] Ingin menyerang atau melarikan diri
    \item[$\square$] Tidak bisa mendengar/menerima informasi
\end{itemize}

\textbf{STAGE 2: STOP (Jika 3+ gejala terdeteksi)}
\begin{enumerate}[itemsep=0.15cm]
    \item[$\square$] Gunakan safe word atau katakan: ``Aku flooded, aku butuh break.''
    \item[$\square$] Tidak ada ``kata terakhir'' --- langsung stop
    \item[$\square$] Beri tahu kapan akan kembali: ``Aku akan kembali dalam 20-30 menit.''
    \item[$\square$] Pergi ke ruang terpisah
\end{enumerate}

\textbf{STAGE 3: SELF-SOOTHING (Selama break)}
\begin{itemize}[itemsep=0.1cm]
    \item[$\square$] Napas dalam 4-7-8 (tarik 4 detik, tahan 7, hembus 8)
    \item[$\square$] Meditasi atau mindfulness
    \item[$\square$] Olahraga ringan (jalan kaki)
    \item[$\square$] Mandi air hangat
    \item[$\square$] Dengar musik menenangkan
    \item[$\square$] \textbf{JANGAN}: Pikirkan argumen, merencanakan balas dendam, stalking medsos
\end{itemize}

\textbf{STAGE 4: RETURN (Setelah 20-30 menit)}
\begin{enumerate}[itemsep=0.15cm]
    \item[$\square$] Pastikan sudah tenang (jantung normal)
    \item[$\square$] Kembali ke pasangan sesuai janji
    \item[$\square$] Mulai dengan soft start-up
    \item[$\square$] Jika masih flooded, minta tambahan waktu dengan spesifik
\end{enumerate}
\end{warningbox}

\subsection{Protokol: Ketika Repair Attempts Gagal}

\begin{warningbox}[EMERGENCY PROTOCOL --- REPAIR ATTEMPTS GAGAL]
\textbf{Jika repair attempts kamu terus ditolak:}

\vspace{0.2cm}
\textbf{Langkah 1: Evaluasi}
\begin{itemize}[itemsep=0.1cm]
    \item[$\square$] Apakah waktunya tepat? (Jika pasangan masih flooded, tunggu)
    \item[$\square$] Apakah cara repair-nya sesuai dengan love language pasangan?
    \item[$\square$] Apakah kamu sudah mengakui bagian yang salah dengan spesifik?
\end{itemize}

\textbf{Langkah 2: Coba Pendekatan Berbeda}
\begin{itemize}[itemsep=0.1cm]
    \item[$\square$] Jika verbal gagal $\rightarrow$ Coba tulis surat/catatan
    \item[$\square$] Jika langsung gagal $\rightarrow$ Coba melalui aktivitas (cooking, jalan bareng)
    \item[$\square$] Jika serius gagal $\rightarrow$ Coba humor ringan atau affection fisik
\end{itemize}

\textbf{Langkah 3: Cek Deeper Issue}
\begin{itemize}[itemsep=0.1cm]
    \item[$\square$] Apakah ada masalah yang belum terselesaikan dari konflik sebelumnya?
    \item[$\square$] Apakah pasangan merasa tidak didengar/dipahami?
    \item[$\square$] Apakah ada trust issue yang mendalam?
\end{itemize}

\textbf{Langkah 4: Escalate}
\begin{itemize}[itemsep=0.1cm]
    \item[$\square$] Jadwalkan State of the Union meeting khusus untuk ini
    \item[$\square$] Pertimbangkan couples therapy jika pola ini berulang
    \item[$\square$] Jangan memaksa repair --- beri waktu, tapi tetap tunjukkan komitmen
\end{itemize}
\end{warningbox}

\subsection{Protokol: Ketika Mempertimbangkan Separation (Last Resort)}

\begin{warningbox}[EMERGENCY PROTOCOL --- LAST RESORT CHECKLIST]
\textbf{Sebelum memutuskan untuk berpisah, pastikan:}

\vspace{0.2cm}
\textbf{CHECKLIST TINDAKAN INTENSIF}

\textit{Sudahkah kalian:}
\begin{itemize}[itemsep=0.15cm]
    \item[$\square$] Melakukan State of the Union meeting rutin selama minimal 3 bulan?
    \item[$\square$] Mencoba couples therapy dengan profesional berlisensi?
    \item[$\square$] Menerapkan semua 7 Principles secara konsisten?
    \item[$\square$] Menerapkan safe word dan flooding protocol?
    \item[$\square$] Mencoba 30-day intensive connection program?
    \item[$\square$] Membaca dan menerapkan materi dari buku-buku rekomendasi?
    \item[$\square$] Memiliki support system (teman, keluarga, komunitas)?
\end{itemize}

\textbf{CHECKLIST EVALUASI HUBUNGAN}

\textit{Apakah masih ada:}
\begin{itemize}[itemsep=0.15cm]
    \item[$\square$] Fondness and admiration (meski kecil)?
    \item[$\square$] Friendship dasar?
    \item[$\square$] Shared meaning (visi bersama)?
    \item[$\square$] Kemauan dari kedua belah pihak untuk berusaha?
    \item[$\square$] Physical safety (tidak ada abuse)?
\end{itemize}

\textbf{JIKA ADA ABUSE (Fisik/Emosional):}
\begin{itemize}[itemsep=0.15cm]
    \item[$\square$] Prioritaskan keselamatan fisik dan emosionalmu
    \item[$\square$] Hubungi hotline darurat: 119 (Indonesia) atau hotline lokal
    \item[$\square$] Cari tempat aman
    \item[$\square$] Abuse bukan alasan untuk ``coba lagi'' --- safety first
\end{itemize}

\textbf{LANGKAH FINAL JIKA MEMUTUSKAN BERPISAH:}
\begin{enumerate}[itemsep=0.15cm]
    \item[$\square$] Konsultasi dengan counselor individual
    \item[$\square$] Siapkan support system
    \item[$\square$] Buat rencana praktis (finansial, tempat tinggal)
    \item[$\square$] Komunikasikan dengan jujur tapi sopan
    \item[$\square$] Terapkan closure ritual untuk berdua
\end{enumerate}
\end{warningbox}

\section{Resource Guide (Panduan Sumber Daya)}

\subsection{Further Reading (Buku yang Direkomendasikan)}

\begin{praktikbox}[Rekomendasi Bacaan]
\textbf{Foundational (Buku Utama):}
\begin{itemize}[itemsep=0.15cm]
    \item \textbf{The Seven Principles for Making Marriage Work} oleh John Gottman --- Buku dasar yang menjadi fondasi panduan ini
    \item \textbf{Wired for Dating} dan \textbf{Wired for Love} oleh Stan Tatkin --- Memahami attachment styles dan neuroscience relationship
    \item \textbf{Hold Me Tight} oleh Sue Johnson --- Emotionally Focused Therapy (EFT) untuk pasangan
    \item \textbf{Attached} oleh Amir Levine \& Rachel Heller --- Attachment theory dalam konteks dating modern
\end{itemize}

\textbf{Advanced Reading:}
\begin{itemize}[itemsep=0.15cm]
    \item \textbf{What Makes Love Last?} oleh John Gottman --- Tentang trust dan betrayal
    \item \textbf{The Relationship Cure} oleh John Gottman --- Emotional connection dalam semua relasi
    \item \textbf{Love Sense} oleh Sue Johnson --- Science of love attachment
    \item \textbf{Eight Dates} oleh John \& Julie Gottman --- Guided conversations untuk pasangan
    \item \textbf{Nonviolent Communication} oleh Marshall Rosenberg --- Komunikasi dengan empati
\end{itemize}

\textbf{For Specific Issues:}
\begin{itemize}[itemsep=0.15cm]
    \item \textbf{Come As You Are} oleh Emily Nagoski --- Sexual health dan intimacy
    \item \textbf{The Five Love Languages} oleh Gary Chapman --- Expressions of love
    \item \textbf{Getting the Love You Want} oleh Harville Hendrix --- Imago Relationship Therapy
    \item \textbf{Mating in Captivity} oleh Esther Perel --- Erotic intelligence dalam long-term relationship
\end{itemize}
\end{praktikbox}

\subsection{Kapan Harus Mencari Bantuan Profesional}

\begin{warningbox}[Tanda-tanda Butuh Terapis]
\textbf{Segera cari bantuan profesional jika:}
\begin{itemize}[itemsep=0.15cm]
    \item[$\square$] Four Horsemen dominan dan repair attempts gagal terus-menerus
    \item[$\square$] Salah satu atau kedua pihak sering flooded dan tidak bisa self-soothe
    \item[$\square$] Adanya kekerasan fisik atau ancaman kekerasan
    \item[$\square$] Adanya perselingkuhan (emotional atau fisik) yang belum terselesaikan
    \item[$\square$] Salah satu pihang mengancam perceraian/berpisah secara teratur
    \item[$\square$] Tidak ada fondness and admiration yang tersisa
    \item[$\square$] Konflik perpetual tidak bisa didialogkan dengan baik
    \item[$\square$] Tidak ada kemajuan setelah 3-6 bulan mencoba sendiri
\end{itemize}

\textbf{Jenis Terapis yang Direkomendasikan:}
\begin{itemize}[itemsep=0.1cm]
    \item \textbf{Gottman Method Therapist} --- Terlatih dalam metode Gottman
    \item \textbf{EFT Therapist} --- Emotionally Focused Therapy (ICEEFT certified)
    \item \textbf{PACT Therapist} --- Psychobiological Approach to Couple Therapy (Stan Tatkin)
    \item \textbf{Imago Therapist} --- Imago Relationship Therapy
\end{itemize}
\end{warningbox}

\subsection{Sumber Daya Online}

\begin{tipbox}[Resources Digital]
\textbf{Apps untuk Pasangan:}
\begin{itemize}[itemsep=0.1cm]
    \item \textbf{Gottman Card Decks} --- Prompts untuk conversation dan intimacy
    \item \textbf{Lasting} --- Marriage counseling app berbasis Gottman
    \item \textbf{Paired} --- Daily questions dan exercises untuk pasangan
    \item \textbf{Relish} --- Relationship coaching app
\end{itemize}

\textbf{Website \& Online Resources:}
\begin{itemize}[itemsep=0.1cm]
    \item \textbf{Gottman.com} --- Artikel, assessment, dan workshop
    \item \textbf{The Gottman Institute Blog} --- Research-based relationship advice
    \item \textbf{StanTatkin.com} --- Attachment and neuroscience of relationships
    \item \textbf{ICEEFT.com} --- Emotionally Focused Therapy resources
    \item \textbf{Psychology Today} --- Directory untuk mencari therapist lokal
\end{itemize}

\textbf{Assessment Online (Gratis):}
\begin{itemize}[itemsep=0.1cm]
    \item Love Languages Quiz (5lovelanguages.com)
    \item Attachment Style Quiz (various online resources)
    \item Gottman Love Lab assessments
\end{itemize}
\end{tipbox}

\vfill

\begin{center}
\begin{tikzpicture}
    \draw[primary, line width=3pt] (0,0) -- (12,0);
\end{tikzpicture}

\vspace{0.5cm}

{\Large\bfseries\color{primary} Selamat Membangun Hubungan yang Sehat!}

\vspace{0.3cm}

{\color{warmgray} ``Hubungan yang baik bukan tentang tidak pernah bertengkar,\\
tapi tentang bertengkar dengan cara yang memperkuat, bukan merusak.''}

\vspace{0.3cm}

{\color{warmgray} ``Tools dalam bab ini adalah kompas --- gunakan setiap saat\\
untuk menavigasi perjalanan cinta kalian.''}

\vspace{0.5cm}

{\Large\color{accent} $\bullet$ \hspace{0.3cm} $\bullet$ \hspace{0.3cm} $\bullet$}
\end{center}
